\section{Urn models}

Urn models are a standard way to study how probability spaces are constructed.
They are especially useful because one physical setup can be described at
several levels. We may record the exact labelled object selected from the urn,
or we may record only a visible type such as colour. We may sample with
replacement or without replacement. We may record order or ignore order. Each
choice leads to a different sample space.

\subsection*{One selection from labelled objects}

Suppose an urn contains five labelled objects:
\[
R_1,R_2,R_3,B_1,B_2,
\]
where the letter records type and the subscript records the label. If one object
is selected and its label is recorded, then
\[
\Omega=\{R_1,R_2,R_3,B_1,B_2\}.
\]
If the selection mechanism treats the five labelled objects equally, then each
has probability \(1/5\).

The event that the selected object is of type \(R\) is
\[
A=\{R_1,R_2,R_3\}.
\]
Therefore
\[
P(A)=\frac35.
\]
The event that the selected object is of type \(B\) is
\[
A^c=\{B_1,B_2\},
\]
so
\[
P(A^c)=\frac25.
\]

\begin{remarkbox}
In this model, \(R_1\), \(R_2\), and \(R_3\) are different elementary outcomes.
They have the same type, but the labels are recorded. The event ``type \(R\)''
is therefore a set of elementary outcomes, not a single elementary outcome.
\end{remarkbox}

\subsection*{Recording only type}

Now use the same physical urn, but record only the type of the selected object.
Then the sample space is
\[
\Omega'=\{R,B\}.
\]
This sample space has two elements, but it is not automatically uniform. The
probability of type \(R\) is inherited from the three labelled objects of type
\(R\), and the probability of type \(B\) is inherited from the two labelled
objects of type \(B\):
\[
P(R)=\frac35,\qquad P(B)=\frac25.
\]

This is the central lesson of the urn model: when several detailed elementary
outcomes are grouped into one recorded outcome, equal probabilities at the
detailed level may become unequal probabilities at the aggregated level.

\begin{examplebox}
Writing
\[
P(R)=P(B)=\frac12
\]
would describe a different mechanism: one that first chooses a type uniformly.
It would not describe uniform selection from an urn containing three labelled
objects of type \(R\) and two labelled objects of type \(B\).
\end{examplebox}

\subsection*{Two selections with replacement}

Suppose one object is selected, recorded, returned to the urn, and then another
selection is made. If labels and order are recorded, then one elementary outcome
is an ordered pair
\[
(x_1,x_2),
\]
where each coordinate belongs to
\[
\{R_1,R_2,R_3,B_1,B_2\}.
\]
Because the first selected object is returned, the same label may appear twice.
Thus
\[
\Omega=\{R_1,R_2,R_3,B_1,B_2\}^2,\qquad |\Omega|=25.
\]

The event that both selected objects have type \(R\) is
\[
A=\{(R_i,R_j):i,j\in\{1,2,3\}\}.
\]
There are \(3\cdot3=9\) such ordered pairs, so
\[
P(A)=\frac9{25}.
\]
Equivalently,
\[
P(A)=\frac35\cdot\frac35.
\]

\subsection*{Two selections without replacement}

Now suppose the first selected object is not returned before the second
selection. If labels and order are recorded, then
\[
\Omega=\{(x_1,x_2):x_1,x_2\in\{R_1,R_2,R_3,B_1,B_2\},\ x_1\neq x_2\}.
\]
There are
\[
5\cdot4=20
\]
elementary outcomes.

The event that both selected objects have type \(R\) is
\[
A=\{(R_i,R_j): i,j\in\{1,2,3\},\ i\neq j\}.
\]
There are \(3\cdot2=6\) such ordered pairs, hence
\[
P(A)=\frac6{20}=\frac3{10}.
\]
The probability is smaller than in the replacement model because, after one
object of type \(R\) has been selected, fewer such objects remain available.

\begin{examplebox}
The event ``one object of each type'' has two ordered forms:
\[
(R_i,B_j)\quad\text{or}\quad(B_j,R_i).
\]
There are \(3\cdot2=6\) outcomes of the first form and \(2\cdot3=6\) outcomes of
the second form. Therefore
\[
P(\text{one of each type})=\frac{12}{20}=\frac35.
\]
\end{examplebox}

\subsection*{Order recorded or ignored}

If two objects are selected without replacement and order is ignored, then an
elementary outcome is a two-element subset of the five labelled objects:
\[
\Omega=\{S\subseteq\{R_1,R_2,R_3,B_1,B_2\}: |S|=2\}.
\]
Thus
\[
|\Omega|=\binom52=10.
\]

The event that both selected objects have type \(R\) has
\[
\binom32=3
\]
favourable subsets, so
\[
P(\text{both type }R)=\frac{\binom32}{\binom52}=\frac3{10}.
\]
This agrees with the ordered calculation because each unordered pair corresponds
to exactly two ordered pairs.

\subsection*{Only the number of type \(R\) objects is recorded}

Sometimes we ignore labels and order and record only a count. In two selections
with replacement, the possible recorded values for the number of type \(R\)
objects are
\[
\Omega'=\{0,1,2\}.
\]
The probabilities are not uniform. Since \(P(R)=3/5\) and \(P(B)=2/5\),
\[
P(0)=\left(\frac25\right)^2=\frac4{25},
\]
\[
P(1)=2\cdot\frac35\cdot\frac25=\frac{12}{25},
\]
\[
P(2)=\left(\frac35\right)^2=\frac9{25}.
\]
The middle value has a factor \(2\) because exactly one type \(R\) object can
appear in two orders: first or second.

\subsection*{General formula with replacement}

Suppose the urn contains \(r\) labelled objects of type \(R\) and \(b\) labelled
objects of type \(B\). If one object is selected uniformly and only type is
recorded, then
\[
P(R)=\frac{r}{r+b},\qquad P(B)=\frac{b}{r+b}.
\]

If \(n\) selections are made with replacement and types are recorded in order,
then
\[
\Omega=\{R,B\}^n.
\]
The probability of exactly \(k\) recorded symbols of type \(R\) is
\[
\binom nk
\left(\frac{r}{r+b}\right)^k
\left(\frac{b}{r+b}\right)^{n-k}.
\]
The binomial coefficient counts the positions occupied by type \(R\).

\subsection*{General formula without replacement}

If \(n\) objects are selected without replacement and order is ignored, then the
total number of possible labelled selections is
\[
\binom{r+b}{n}.
\]
To obtain exactly \(k\) objects of type \(R\), choose \(k\) from the \(r\) objects
of type \(R\), and choose \(n-k\) from the \(b\) objects of type \(B\). Thus
\[
P(\text{exactly }k\text{ of type }R)
=
\frac{\binom rk\binom b{n-k}}{\binom{r+b}{n}}.
\]

\begin{examplebox}
If \(r=7\), \(b=5\), and \(n=4\), then the probability of selecting exactly two
objects of type \(R\), without replacement and ignoring order, is
\[
\frac{\binom72\binom52}{\binom{12}{4}}.
\]
\end{examplebox}

\subsection*{What urn models teach}

Urn models show that probability depends on both the mechanism and the recorded
information. A detailed uniform model may become non-uniform after aggregation.
Replacement keeps the composition of the urn fixed from selection to selection;
without replacement, the composition changes. Recording order gives sequences;
ignoring order gives sets or count summaries.

\begin{remarkbox}
The phrase ``select from an urn'' is not yet a probability model. A complete
model must say what is recorded, whether replacement occurs, whether order is
recorded, and whether labelled objects or only their types are treated as
elementary outcomes.
\end{remarkbox}
